\section{Limitations}

\textsc{ProprioVLA} assumes that the robot is at least partially visible during
training. If the embodiment is consistently outside the camera view, the
state-aligned motion signal cannot provide useful self-grounding supervision,
and the visual query has no direct evidence from which to recover the robot's
own state.

State-aligned residual motion can be ambiguous when non-robot motion is coupled
to the robot. Grasped objects, deformable materials, moving backgrounds, or
other agents may move synchronously with proprioceptive change and receive high
alignment scores. Confidence weighting reduces weak supervision, but it cannot
fully distinguish the robot body from every object moved by that body.

Finally, the evaluation covers a finite set of embodiments, camera
perturbations, and manipulation tasks. Broader multi-camera setups, mobile
manipulation, multi-robot scenes, long-horizon tasks, and highly dynamic
environments remain important tests for whether proprioceptive self-grounding
scales beyond controlled perturbation regimes.
